\unnumberedsection{Заключение}

В ходе выполнения дипломного проекта было разработано веб-приложение для управления приёмом абитуриентов в Белорусский национальный технический университет.
В процессе работы проведён анализ предметной области, исследованы существующие программные аналоги, выделены их ключевые ограничения, сформулированы функциональные, технические и бизнес-требования к разрабатываемой системе.

Разработка велась с использованием современных технологий ASP.NET~Core, React, PostgreSQL и Liquibase.
Клиент-серверная архитектура, реализованная в проекте, обеспечила структурированность, гибкость и масштабируемость приложения, а также возможность независимого развития серверной и клиентской частей.

Разработанная система обеспечивает:
\begin{itemize}
\item ведение организационной структуры университета (факультеты, кафедры, специальности);
\item ведение конкурсных списков и категорий приёма с указанием плана и квот;
\item регистрацию абитуриентов и приём их заявлений в категории приёма с указанием приоритета предпочтения;
\item автоматический пересчёт результатов отбора с учётом квот категорий и общего плана конкурсного списка;
\item экспорт результатов в формат Excel для передачи в государственные информационные системы;
\item разграничение прав доступа на основе ролевой модели;
\item журналирование всех действий пользователей и двойной контроль критичных операций.
\end{itemize}

Проведённое функциональное тестирование подтвердило корректность работы всех реализованных функций и соответствие приложения сформулированным требованиям.
Алгоритм пересчёта результатов отбора продемонстрировал корректное поведение на всех проверенных сценариях, включая граничные случаи с превышением квот категорий и общего плана конкурсного списка.
Дополнительно была выполнена валидация на реальных данных приёмной кампании БНТУ за 2024~год: результаты, рассчитанные системой, полностью совпали с фактическими итогами приёма как по персональному составу зачисленных абитуриентов, так и по проходным баллам по всем конкурсным спискам, что подтверждает применимость разработанного решения в условиях реальной приёмной кампании.

Результаты проекта могут быть применены в работе приёмных комиссий учреждений высшего образования Республики Беларусь.
Дальнейшее развитие системы может включать интеграцию с государственными информационными системами для автоматизированного получения сведений о результатах централизованного тестирования абитуриентов, добавление модуля онлайн-подачи заявлений непосредственно абитуриентами, реализацию мобильного клиента для контроля хода приёмной кампании руководством университета.

Таким образом, поставленная цель дипломного проекта достигнута, все сформулированные задачи выполнены, а полученные результаты обладают практической ценностью и могут служить основой для дальнейшего развития и масштабирования системы.\label{lastmainpage}
